
\subsubsection{5.1 Existence of Semantic Accommodation (h-e1)}

Table 1 reports the partner-specificity hierarchy for the primary model (all-MiniLM-L6-v2, n = 155,362 pairs).

\textbf{Table 1: Partner-Specificity Hierarchy (all-MiniLM-L6-v2)}

\begin{table}[htbp]
\centering
\resizebox{\columnwidth}{!}{%
\begin{tabular}{llll}
\toprule
Control Level & Mean Cosine Similarity & 95\% CI & Cohen's d vs. Actual \\
\midrule
Actual AI partner & 0.3534 & [0.352, 0.354] & -- \\
KNN topic-matched (K=5) & 0.2688 & [0.268, 0.270] & 0.417 \\
Random AI turn & 0.0241 & [0.024, 0.025] & 1.998 \\
\bottomrule
\end{tabular}
}
\end{table}

C\_sem = 0.3534 - 0.0241 = 0.329 (95\% CI [0.328, 0.330]). The ordering actual > KNN > random holds, with all Mann-Whitney comparisons significant at p < 0.001 (p = 0.0 at machine precision). The effect size relative to the random baseline (d = 1.998) is large by conventional standards \citep{Cohen1988Statistical}. The effect relative to the KNN topic-matched baseline (d = 0.417) is smaller but remains above zero, indicating that the signal is not fully explained by topical co-occurrence.

Robustness models confirm the pattern. For paraphrase-MiniLM-L6-v2: C\_sem = 0.300 (95\% CI [0.299, 0.301]), d = 1.838 vs. random, d = 0.413 vs. KNN. For all-mpnet-base-v2: C\_sem = 0.341 (95\% CI [0.340, 0.342]), d = 2.099 vs. random, d = 0.400 vs. KNN. All three models pass the existence gate.

\begin{figure}[htbp]
\centering
\includegraphics[width=\columnwidth]{figures/partner_specificity.png}
\caption{Partner-specificity hierarchy}
\label{fig:partner_specificity}
\end{figure}

\textbf{Figure 1.} Three-level partner-specificity hierarchy showing mean cosine similarity for actual AI partner, KNN topic-matched (K=5), and random AI turns (all-MiniLM-L6-v2, n = 155,362).

\subsubsection{5.2 Tier-Scalable Accommodation (h-m1)}

Table 2 reports C\_sem values across tiers and models.

\textbf{Table 2: C\_sem per Tier per SBERT Model}

\begin{table}[htbp]
\centering
\resizebox{\columnwidth}{!}{%
\begin{tabular}{llll}
\toprule
Tier & all-MiniLM-L6-v2 & paraphrase-MiniLM-L6-v2 & all-mpnet-base-v2 \\
\midrule
T1 (helpful-base) & 0.3036 & 0.2714 & 0.3138 \\
T2 (helpful-rejection-sampled) & 0.3367 & 0.3068 & 0.3483 \\
T3 (helpful-online) & 0.3678 & 0.3456 & 0.3820 \\
\bottomrule
\end{tabular}
}
\end{table}

Monotonicity (T1 < T2 < T3) holds for all three models. IPW-corrected aggregate values (averaged across models) are 0.307, 0.336, and 0.364 for T1, T2, and T3 respectively; monotonicity is preserved after IPW correction.

\textbf{Table 3: Jonckheere-Terpstra Tests and Effect Sizes (h-m1)}

\begin{table*}[htbp]
\centering
\resizebox{\dimexpr 2\columnwidth+\columnsep\relax}{!}{%
\begin{tabular}{llllll}
\toprule
Model & J-T Statistic & J-T p-value & d (T1 vs. T2) & d (T2 vs. T3) & d (T1 vs. T3) \\
\midrule
all-MiniLM-L6-v2 & 4,039,609,994 & 0.001 & 0.087 & 0.096 & 0.183 \\
paraphrase-MiniLM-L6-v2 & 4,116,983,177 & 0.001 & 0.114 & 0.143 & 0.254 \\
all-mpnet-base-v2 & 4,090,937,042 & 0.001 & 0.098 & 0.140 & 0.238 \\
\bottomrule
\end{tabular}
}
\end{table*}

All three models yield J-T p = 0.001. The T1-to-T3 Cohen's d values range from 0.183 to 0.254, all exceeding the pre-specified threshold of 0.1. Adjacent-tier contrasts (T1 vs. T2) yield smaller effect sizes (d = 0.087--0.114), with some falling below the 0.1 threshold. The Jonckheere-Terpstra test evaluates the full ordered pattern rather than adjacent contrasts alone.

\subsubsection{5.3 Directional Asymmetry (h-m2)}

Table 4 reports C\_sem in both directions across all nine tier-by-model conditions.

\textbf{Table 4: Directional C\_sem Asymmetry}

\begin{table*}[htbp]
\centering
\resizebox{\dimexpr 2\columnwidth+\columnsep\relax}{!}{%
\begin{tabular}{llllll}
\toprule
Model & Tier & $C_sem^{H<-A}$ & $C_sem^{A<-H}$ & Cohen's d & p-value \\
\midrule
all-MiniLM-L6-v2 & T1 (base) & 0.0853 & 0.0395 & 0.373 & 0.0 \\
all-MiniLM-L6-v2 & T2 (rej.-sampled) & 0.0923 & 0.0535 & 0.330 & 0.0 \\
all-MiniLM-L6-v2 & T3 (online) & 0.0876 & 0.0718 & 0.133 & 4.8e-30 \\
paraphrase-MiniLM-L6-v2 & T1 & 0.0794 & 0.0316 & 0.405 & 0.0 \\
paraphrase-MiniLM-L6-v2 & T2 & 0.0866 & 0.0433 & 0.385 & 0.0 \\
paraphrase-MiniLM-L6-v2 & T3 & 0.0847 & 0.0617 & 0.205 & 4.5e-81 \\
all-mpnet-base-v2 & T1 & 0.0826 & 0.0422 & 0.334 & 0.0 \\
all-mpnet-base-v2 & T2 & 0.0884 & 0.0581 & 0.261 & 0.0 \\
all-mpnet-base-v2 & T3 & 0.0838 & 0.0767 & 0.061 & 0.004 \\
\bottomrule
\end{tabular}
}
\end{table*}

$C_sem^{H<-A}$ exceeds $C_sem^{A<-H}$ in all nine conditions. Effect sizes range from d = 0.061 (mpnet, T3) to d = 0.405 (paraphrase, T1). All comparisons are statistically significant. The weakest cell (all-mpnet-base-v2, T3 online) has d = 0.061 with p = 0.004. The asymmetry tends to be smaller for T3 (the highest-quality tier) across all models.

\begin{figure}[htbp]
\centering
\includegraphics[width=\columnwidth]{figures/significance_heatmap.png}
\caption{Directional asymmetry}
\label{fig:significance_heatmap}
\end{figure}

\textbf{Figure 2.} Cohen's d heatmap for directional asymmetry ($C_sem^{H<-A}$ minus $C_sem^{A<-H}$) across all nine tier-by-model cells.

\subsubsection{5.4 Within-Prompt Quality Discrimination (h-m3)}

Table 5 reports delta = cos(H\_next, A\_chosen) - cos(H\_next, A\_rejected) across tiers and operationalizations for all-MiniLM-L6-v2.

\textbf{Table 5: Within-Prompt Delta (all-MiniLM-L6-v2)}

\begin{table*}[htbp]
\centering
\resizebox{\dimexpr 2\columnwidth+\columnsep\relax}{!}{%
\begin{tabular}{llllllll}
\toprule
Tier & n pairs & OP1 (raw) E[delta] & OP1 d & OP2 (length-matched) E[delta] & OP2 d & OP3 (prompt-projected) E[delta] & OP3 d \\
\midrule
T1 (base) & 31,013 & -0.020 & -0.074 & -0.041 & -0.161 & +0.014 & +0.069 \\
T2 (rej.-sampled) & 35,665 & -0.106 & -0.400 & -0.118 & -0.459 & -0.045 & -0.216 \\
T3 (online) & 14,426 & -0.167 & -0.716 & -0.170 & -0.738 & -0.097 & -0.524 \\
\bottomrule
\end{tabular}
}
\end{table*}

Delta is negative in 25 of 27 tier-by-operationalization-by-model combinations across all three SBERT models. The single partial exception is OP3 (prompt-projected) for T1 with all-MiniLM-L6-v2, which yields a weakly positive delta of +0.014 (d = +0.069). Zero of three models pass the pre-specified gate (delta > 0 in at least two of three operationalizations).

The negative delta is strongest in T3 (d = -0.716 to -0.738 for OP1 and OP2), the tier with the highest RLHF quality. Human follow-up turns are systematically more semantically similar to the rejected AI response than to the chosen one. This pattern is statistically significant at p < 0.001 in 24 of 27 cells.

The paraphrase-MiniLM-L6-v2 and all-mpnet-base-v2 models show qualitatively identical patterns, with all operationalizations yielding negative deltas except for isolated weak positive values in T1 prompt-projected conditions. Neither robustness model passes the gate.

\subsubsection{5.5 Politeness-Style Mediation (h-m4)}

Table 6 reports OLS mediation regression results.

\textbf{Table 6: OLS Mediation Regression (n = 3,000 per model)}

\begin{table*}[htbp]
\centering
\resizebox{\dimexpr 2\columnwidth+\columnsep\relax}{!}{%
\begin{tabular}{lllll}
\toprule
Model & beta\_PM & SE (HC3) & p-value & R-squared \\
\midrule
all-MiniLM-L6-v2 & -1.46e-05 & ~1.5e-02 & 0.998 & 0.007 \\
paraphrase-MiniLM-L6-v2 & -1.26e-06 & ~1.2e-02 & 1.000 & 0.007 \\
all-mpnet-base-v2 & +6.76e-05 & ~7.2e-03 & 0.991 & 0.012 \\
\bottomrule
\end{tabular}
}
\end{table*}

beta\_PM is indistinguishable from zero in all three models (absolute values below 1e-4, p-values near 1.0). Total model R-squared ranges from 0.007 to 0.012, indicating that the full set of predictors (PM-proxy, surface features, tier controls) collectively explains less than 1.2\% of C\_sem variance. Zero of three models pass the mediation gate.

\subsubsection{5.6 Summary}

\textbf{Table 7: Sub-Hypothesis Verification Summary}

\begin{table}[htbp]
\centering
\resizebox{\columnwidth}{!}{%
\begin{tabular}{llll}
\toprule
Sub-Hypothesis & Gate Level & Result & Key Metric \\
\midrule
h-e1: Accommodation exists & MUST\_WORK & PASS & C\_sem = 0.329; d = 1.998 vs. random \\
h-m1: Tier monotonicity & MUST\_WORK & PASS & J-T p = 0.001; d = 0.183--0.254 (T1 to T3); 3/3 models \\
h-m2: Directional asymmetry & SHOULD\_WORK & PASS & All 9 cells; d = 0.061--0.405 \\
h-m3: Within-prompt quality discrimination & SHOULD\_WORK & FAIL & delta < 0 in 25/27 cells; d up to -0.738 \\
h-m4: PM-proxy mediation & SHOULD\_WORK & FAIL & beta\_PM approximately 0; p approximately 0.99 \\
\bottomrule
\end{tabular}
}
\end{table}

Three of five sub-hypotheses are supported. Two mechanism hypotheses are falsified.
